% Standalone problem document, generated from the adjacent README.md.
% Compile with XeLaTeX. Mathematical content is maintained in Markdown.
\documentclass[11pt,a4paper]{article}
\usepackage[fontsize=10.5pt]{scrextend}
\usepackage[margin=22mm,headheight=15pt,headsep=9mm,footskip=11mm]{geometry}
\usepackage{fontspec}
\setmainfont{texgyrepagella}[Extension=.otf,UprightFont=*-regular,BoldFont=*-bold,ItalicFont=*-italic,BoldItalicFont=*-bolditalic]
\setsansfont{texgyreheros}[Extension=.otf,UprightFont=*-regular,BoldFont=*-bold,ItalicFont=*-italic,BoldItalicFont=*-bolditalic]
\setmonofont{texgyrecursor}[Extension=.otf,UprightFont=*-regular,BoldFont=*-bold,ItalicFont=*-italic,BoldItalicFont=*-bolditalic,Scale=MatchLowercase]
\usepackage{amsmath,amssymb}
\usepackage{unicode-math}
\setmathfont{texgyrepagella-math.otf}
\usepackage{xcolor,microtype,enumitem,needspace,fancyhdr}
\usepackage{bookmark}
\definecolor{ink}{HTML}{17344B}
\definecolor{accent}{HTML}{197D80}
\definecolor{muted}{HTML}{596773}
\hypersetup{colorlinks=true,linkcolor=accent,urlcolor=accent,pdftitle={TR-17 - Frobenius
inner products minimize the algebraic complexity of rank-one
approximation},pdfauthor={Open Problems in Numerical Linear Algebra}}
\urlstyle{same}
\setlength{\parindent}{0pt}
\setlength{\parskip}{5pt plus 1pt minus 1pt}
\linespread{1.04}
\setlength{\emergencystretch}{3em}
\widowpenalty=10000
\clubpenalty=10000
\displaywidowpenalty=10000
\allowdisplaybreaks[1]
\setlist{nosep,leftmargin=1.5em}
\providecommand{\tightlist}{\setlength{\itemsep}{0pt}\setlength{\parskip}{0pt}}
\providecommand{\pandocbounded}[1]{#1}
\pagestyle{fancy}
\fancyhf{}
\fancyhead[L]{\sffamily\footnotesize\color{muted}OPEN PROBLEMS IN NUMERICAL LINEAR ALGEBRA}
\fancyhead[R]{\sffamily\bfseries\footnotesize\color{accent}TR-17}
\fancyfoot[L]{\parbox[b]{0.9\textwidth}{\sffamily\fontsize{8}{10}\selectfont\color{muted}Editorial ratings; a bounded search cannot certify that a problem remains unsolved.\\Literature check: 2026-09-11}}
\fancyfoot[R]{\sffamily\footnotesize\color{muted}\thepage}
\renewcommand{\headrulewidth}{0.3pt}
\renewcommand{\headrule}{\hbox to\headwidth{\color{accent}\leaders\hrule height \headrulewidth\hfill}}
\makeatletter
\renewcommand{\section}{\Needspace{5\baselineskip}\@startsection{section}{1}{0pt}{12pt plus 2pt minus 2pt}{4pt}{\sffamily\bfseries\large\color{ink}}}
\renewcommand{\subsection}{\Needspace{4\baselineskip}\@startsection{subsection}{2}{0pt}{9pt plus 1pt minus 1pt}{3pt}{\sffamily\bfseries\normalsize\color{ink}}}
\makeatother
\setcounter{secnumdepth}{0}
\begin{document}
{\sffamily\small\color{muted}Tensor computations\par}
\vspace{3pt}
{\raggedright\hyphenpenalty=10000\exhyphenpenalty=10000\sffamily\bfseries\fontsize{21}{24}\selectfont\color{ink}Frobenius
inner products minimize the algebraic complexity of rank-one
approximation\par}
\vspace{7pt}
{\sffamily\small\textbf{Difficulty:} challenging\quad\textbf{Importance:} interesting
to the community\par}
{\sffamily\small\bfseries\color{accent}Status: Solved\par}
\vspace{4pt}
{\color{accent}\hrule height 0.8pt}
\vspace{6pt}
\textbf{Rating rationale:} Challenging because local minimality of ED
degree must be strengthened to a global comparison across all positive
definite metrics; community importance concerns weighted approximation
and partially symmetric tensor models.

\textbf{Affirmative resolution recorded 2026-09-11.} Matthew J.
Colbrook's
\href{https://github.com/ajt60gaibb/OpenProblemsInNLA/blob/main/references/colbrook-tensor-metrics-rank-2026-09-11/manuscripts/tr17_solution.pdf}{complete
manuscript, Theorem 1}
(\href{https://github.com/ajt60gaibb/OpenProblemsInNLA/blob/main/references/colbrook-tensor-metrics-rank-2026-09-11/manuscripts/tr17_solution.tex}{LaTeX
source}) proves the stated global inequality for every permitted
Segre--Veronese format and every positive definite real symmetric
bilinear form, with the complex-bilinear critical-point and multiplicity
convention below. The proof gives nonnegative weights for a
Chern--Schwartz--MacPherson coefficient comparison and proves the needed
signed Euler-characteristic positivity even for singular or nonreduced
quadric sections. No genericity restriction on the metric is imposed.
The full original source passed an
\href{https://github.com/ajt60gaibb/OpenProblemsInNLA/blob/main/references/colbrook-tensor-metrics-rank-2026-09-11/verification/reviews/TR-17-review.md}{independent
mathematical agent review}.
\href{https://github.com/ajt60gaibb/OpenProblemsInNLA/blob/main/references/colbrook-tensor-metrics-rank-2026-09-11/README.md}{Submission
and verification record}. This is documented AI-assisted verification,
not external human peer review; no priority claim is made. The original
target and former difficulty and importance ratings are retained below
as historical context.

\subsection{Statement}\label{statement}

Fix integers \(k\ge1\), \(n_j\ge2\), \(d_j\ge1\), with
\(\sum_jd_j\ge3\). Set \[
W_{\mathbb R}=\bigotimes_{j=1}^k\operatorname{Sym}^{d_j}(\mathbb R^{n_j}),\qquad
X=\{c\,v_1^{\otimes d_1}\otimes\cdots\otimes v_k^{\otimes d_k}:c\in\mathbb C,\ v_j\in\mathbb C^{n_j}\}\subset W_{\mathbb C}.
\] For a positive definite symmetric real bilinear form \(Q\) on
\(W_{\mathbb R}\), extend it complex-bilinearly. Define
\(\operatorname{ED}_Q(X)\) to be the number of complex critical points
on the smooth nonzero part of \(X\) of \(Z\mapsto Q(A-Z,A-Z)\), for
Zariski-generic \(A\in W_{\mathbb C}\), with algebraic multiplicities.

Let \(Q_F\) be the restriction of the entrywise Frobenius inner product
from the full tensor product, with standard Euclidean inner products on
the factors. Is \[
\operatorname{ED}_Q(X)\ge\operatorname{ED}_{Q_F}(X)
\] true for every such \(Q\)? The bilinear complexification uses no
complex conjugates; the Frobenius restriction includes the usual
repeated-entry weights on symmetric tensors.

\subsection{Relevance}\label{relevance}

The ED degree counts algebraic stationary solutions of weighted rank-one
approximation. The conjecture says the natural tensor metric minimizes
this complexity, across arbitrary positive definite weights and partial
symmetries.

\subsection{References}\label{references}

\begin{enumerate}
\def\labelenumi{\arabic{enumi}.}
\tightlist
\item
  K. Kozhasov, A. Muniz, Y. Qi, and L. Sodomaco, \emph{On the minimal
  algebraic complexity of the rank-one approximation problem for general
  inner products}. \href{https://arxiv.org/html/2309.15105v3}{Primary
  text}, Conjecture 3.9; Theorem 3.12 and §4;
  \href{https://doi.org/10.1090/mcom/4176}{publication DOI}.
\item
  J. Draisma, E. Horobeţ, G. Ottaviani, B. Sturmfels, and R. R. Thomas,
  \emph{The Euclidean distance degree of an algebraic variety}, Found.
  Comput. Math. (2016). \href{https://arxiv.org/pdf/1309.0049}{Primary
  preprint}, §§1--2 for ED degree and §4 for average ED degree.
\end{enumerate}

\subsection{Status check --- 2026-09-10}\label{status-check-2026-09-10}

Rechecked \href{https://arxiv.org/html/2309.15105v3}{Kozhasov et al.~v3,
Conjecture 3.9, Theorem 3.12 and §5}, and searched for a later general
proof. Global minimality is proved for binary symmetric tensors and
ternary symmetric cubics, which occur in this target; the matrix theorem
concerns excluded order two. Local minimality is known generally but
does not prove the requested global inequality. No full resolution was
located.
\end{document}
